\section{A.7. Real elliptic curves and integral involutions}
\subsection{The integral algebra, including a further correction to R5}

\begin{proposition}[Proved: the categorical statement of 67/091]\label{A7-gluing-modules}
Let $R$ be a ring, not necessarily commutative, and use left $R$-modules. The category of pairs $(M,\sigma)$ with $M$ an $R$-module, multiplication by $2$ injective on $M$, and $\sigma^2=1$, is equivalent to the category of triples $(P,Q,m)$, where $P,Q$ are $2$-torsion-free $R$-modules and
\[
 m\subset(P/2P)\oplus(Q/2Q),\qquad
 m\cap(P/2P)=m\cap(Q/2Q)=0.
\]
Morphisms of triples are pairs of module maps preserving $m$ after reduction modulo $2$. No Noetherian or finite-generation hypothesis is required. In the proof, localization at $2$ is valid because the integer $2$ is central; equivalently, $M[1/2]=M\otimes_{\Z}\Z[1/2]$.
\end{proposition}
\begin{proof}
Given $(M,\sigma)$, set $P=\ker(\sigma-1)$ and $Q=\ker(\sigma+1)$. Since $2$ is injective, $P\cap Q=0$, while
\[
 2x=(x+\sigma x)+(x-\sigma x)
\]
shows $2M\subset P\oplus Q\subset M$. Embed $M$ in $(P\oplus Q)[1/2]$ using the decomposition of $2x$. Then
\[
 m=2M/(2P\oplus2Q)\subset(P/2P)\oplus(Q/2Q).
\]
If $2x\in P$, then $2\sigma x=2x$, whence $x\in P$; therefore the intersection with the first summand is zero. The other summand is identical.

Conversely, define
\[
 M=\{(p/2,q/2)\in(P\oplus Q)[1/2]:
                (\bar{p},\bar{q})\in m\},
 \qquad \sigma(p/2,q/2)=(p/2,-q/2).
\]
This module is $2$-torsion-free, contains $P\oplus Q$, and has the given two eigenspaces. Indeed an element with second component zero has $\bar{p}=0$ by the intersection condition, so it belongs to $P$, and similarly for $Q$. Reduction of $2M$ recovers $m$. Every map commuting with $\sigma$ restricts to maps of eigenspaces, and the formula for $M$ extends every map of triples uniquely. Hence these functors are mutually inverse, naturally on objects and arrows.
\end{proof}

\begin{theorem}[Refuted: connectedness does not repair the line-bundle corollary]\label{A7-domain-counterexample}
The assertion following 67/091 that the two eigenspaces of a nontrivial involution on a rank-two projective module are invertible is false even over a $2$-torsion-free domain. Thus the supplied correction ``needs $\operatorname{Spec}R$ connected (e.g. a domain)'' is insufficient.
A counterexample is
\[
 R=\Z[x,y,z]/(x^2+yz-1),\qquad M=R^2,
 \qquad \sigma=\begin{pmatrix}x&y\\z&-x\end{pmatrix}.
\]
The categorical statement in Proposition~\ref{A7-gluing-modules} remains valid. Its eigenspaces are invertible under either of the sufficient hypotheses: $R$ is a Dedekind domain and $M$ is finite projective of rank two; or $2$ is invertible in $R$, $\operatorname{Spec}R$ is connected, and $M$ is finite projective of rank two. In these two cases we require $\sigma\ne\pm1$.
\end{theorem}
\begin{proof}
Write $S=\Z[x,y,z]$ and $f=x^2+yz-1$. As a polynomial in $y$ over the UFD $\Z[x,z]$, $f$ is primitive and irreducible over its fraction field; hence $f$ is irreducible in $S$ and $R$ is a domain. In particular $2$ is injective. Direct multiplication gives $\sigma^2=1$ over $R$, and its off-diagonal entries show $\sigma\ne\pm1$.

Let $P=\ker(\sigma-1)$. We claim
\[
 P=\operatorname{im}(\sigma+1),\qquad
 \ker(\sigma+1)=\operatorname{im}(\sigma-1).
\]
For the first identity, lift $v\in P$ to $S^2$. There exists $w\in S^2$ with $(\sigma-1)v=fw$, where the same matrix denotes its lift over $S$. Multiplication by $\sigma+1$ gives $fv=f(\sigma+1)w$. Cancellation of the non-zero-divisor $f$ proves $v=(\sigma+1)w$; the other identity is the same argument. Consequently
\[
 P\simeq\operatorname{coker}(\sigma-1:R^2\to R^2).
\]
At the maximal ideal $\mathfrak m=(2,x-1,y,z)$ every entry of $\sigma-1$ lies in $\mathfrak m$. Right exactness of tensor product therefore gives
$P/\mathfrak mP\simeq\mathbf F_2^2$. On the fraction field the nonzero matrix $\sigma+1$ has rank one, so $P$ has rank one. An invertible module localized at $\mathfrak m$ would have a one-dimensional residue fibre, a contradiction.

For the corrected sufficient hypotheses, over a Dedekind domain the eigenspaces are finite torsion-free modules, hence projective; over the fraction field the involution is diagonalizable with both eigenvalues because it is not scalar, so they have rank one. If $2$ is invertible, $(1\pm\sigma)/2$ are complementary idempotents, making the eigenspaces finite projective. Their ranks are locally constant; connectedness and exclusion of $\pm1$ force each rank to be one. These arguments justify the hypotheses in the places where we use the rank-one conclusion below, namely $R=\Z$ and the specified PID.
\end{proof}

\begin{proposition}[Proved: the exact count in 67/100]\label{A7-PID-count}
Suppose $R$ is a PID, $2\ne0$, and $\kappa=R/2R$ is a field. Conjugacy classes of involutions $\sigma\in\operatorname{GL}_2(R)$ with $\sigma\ne\pm1$ consist of one split class and a set of nonsplit classes canonically indexed by
\[
 \kappa^\times/\operatorname{im}(R^\times\to\kappa^\times).
\]
For $R=\Z$ there are two classes in total. For
$R=\Z[(1+\sqrt{-11})/2]$ there are four classes in total, three of them nonsplit. The struck surjectivity proviso on 67/100 is therefore necessary for the assertion that there is just one nonsplit class.
\end{proposition}
\begin{proof}
The two eigenspaces are free of rank one by Theorem~\ref{A7-domain-counterexample}. In Proposition~\ref{A7-gluing-modules}, $m$ is either zero or the graph of a nonzero map between the one-dimensional $\kappa$-spaces $P/2P,Q/2Q$. After choosing generators it is a line of slope $c\in\kappa^\times$. Changing the two generators multiplies $c$ by a quotient of reductions of units of $R$, precisely the indicated subgroup. Conversely, every such line constructs a finite torsion-free rank-two module $M$, which is free over a PID, and two constructed involutions are conjugate exactly when the triples are isomorphic. This proves completeness and disjointness of the list. The split class has matrix $\operatorname{diag}(1,-1)$; over $\Z$ a nonsplit representative is $\left(\begin{smallmatrix}1&1\\0&-1\end{smallmatrix}\right)$.

Put $\omega=(1+\sqrt{-11})/2$. The norm is
\[
 N(m+n\omega)=m^2+mn+3n^2=(m+n/2)^2+11n^2/4.
\]
This ring is norm Euclidean: for any complex number first choose $n$ so that the imaginary difference from $n\omega$ has modulus at most $\sqrt{11}/4$, then choose $m$ so that the real difference from $m+n\omega$ has modulus at most $1/2$. The squared distance is at most $11/16+1/4=15/16<1$. Applied to the quotient of any two elements of the number field, this is Euclidean division and proves that $R$ is a PID. The norm formula shows that its units are exactly $\pm1$. Finally $\omega$ satisfies $\omega^2-\omega+3=0$, which reduces to the irreducible polynomial $t^2+t+1$ over $\mathbf F_2$. Thus $\kappa=\mathbf F_4$, and the unit image is trivial while $\kappa^\times$ has order three. Explicit representatives of the three nonsplit classes are
\[
 \begin{pmatrix}1&c\\0&-1\end{pmatrix},
 \qquad c\in\{1,\omega,1+\omega\}.
\]
Indeed their eigenlattices are generated by $(1,0)$ and $(-c,2)$, and the gluing line is spanned by $(\bar{c},1)$; the three nonzero residues give the three unit orbits. This proves the count and verifies every hypothesis of the counterexample. The finite residue-field and norm checks are included in \texttt{part\_a567\_checks.py}.
\end{proof}

\subsection{The real analytic and algebraic classification}
We use complex uniformization in its exact form: every complex elliptic curve is $\C/\Lambda$ for a lattice $\Lambda$, and origin-preserving isomorphisms are induced by complex-linear maps taking one lattice onto the other. This is \cite[III, Theorem 3.10 and Proposition 3.3]{MilneEC2021A567}; the relevant statements were opened. The Weierstrass functions satisfy $\wp_{c\Lambda}(cz)=c^{-2}\wp_\Lambda(z)$ and $\wp_{c\Lambda}^{\prime}(cz)=c^{-3}\wp_\Lambda^{\prime}(z)$ (III, \S2 and Proposition 3.7). Thus the linear isomorphisms used here are algebraic in the Weierstrass models. For a conjugation-stable lattice the defining Eisenstein series are real, giving the asserted real descent directly.

\begin{theorem}[Known with complete derivation: classification on 67/093--099]\label{A7-real-classification}
For $a>0$ define real elliptic curves by their complex uniformizations and the indicated conjugation:
\[
 E_a:\quad \C/(\Z+i\sqrt a\,\Z),\qquad
 E'_a:\quad \C/\bigl(\Z+\tfrac{1+i\sqrt a}{2}\Z\bigr),
 \qquad \sigma[z]=[\bar{z}].
\]
These two disjoint families, with the parameter $a$ unique in each, are the complete list of origin-marked real elliptic curves up to real isomorphism. Their real origin-preserving automorphism groups are $\{\pm1\}$.

A fixed complex elliptic curve has real forms if and only if its usual modular invariant $j$ is real, and then has exactly two real isomorphism classes. Outside the square case they are either $\{E_a,E_{1/a}\}$ or $\{E'_a,E'_{1/a}\}$ with $a>1$. For the square lattice they are $\{E_1,E'_1\}$. For the hexagonal lattice they are $\{E'_3,E'_{1/3}\}$.
\end{theorem}
\begin{proof}
A real structure fixing the origin lifts to $z\mapsto\alpha\bar{z}$ on $\C$, preserving $\Lambda$, with $|\alpha|=1$. A complex-linear change of coordinate makes it conjugation. Its invariant and anti-invariant sublattices have generators $u>0$ and $v\in i\R_{>0}$. Applying Proposition~\ref{A7-gluing-modules} over $\Z$, the lattice is either $\Z u\oplus\Z v$ or this lattice enlarged by $(u+v)/2$. Rescale $u$ to $1$, and write $v=i\sqrt a$. These are exactly the two displayed lattices. Conversely, both lattices are conjugation-stable and hence define real curves by descent.

A complex-linear isomorphism commuting with conjugation is multiplication by a nonzero real scalar. It takes the two eigensublattices to the two eigensublattices, preserving their index in the full lattice and their squared-length ratio $a$. This proves uniqueness of the family and parameter. For an automorphism of one normalized lattice the scalar takes the positive invariant generator to itself or its negative, so it is $1$ or $-1$.

For completeness consider all real structures on a fixed complex curve. The group $G^+$ of complex-linear lattice automorphisms is finite cyclic: area preservation gives modulus one, and discreteness of the lattice gives finiteness. The trace of such an automorphism on $\Lambda$ is $2\operatorname{Re}\alpha\in\Z$. Thus the possible rotation orders are $1,2,3,4,6$; since $-1$ is always present, $|G^+|$ is $2,4$, or $6$. If an anti-linear lattice automorphism exists, it is a reflection, and together these maps form the dihedral group of order $2|G^+|$. Since $|G^+|$ is even, conjugation by rotations has exactly two orbits on the reflections: writing them as $r^m\sigma$, conjugation changes $m$ by an even integer.

The existence of a reflection is equivalent to $E\simeq\overline{E}$. The classification by $j$ \cite[II, Proposition 2.1(c)]{MilneEC2021A567} makes this equivalent to $j\in\R$. An anti-linear lattice isomorphism witnessing $E\simeq\overline{E}$ has modulus one by area preservation, hence squares to the identity; so there is no further descent obstruction in this origin-marked setting.

Replacing $\sigma$ by $-\sigma$ interchanges its two eigenlattices, takes $a$ to $1/a$, and leaves their index unchanged. It interchanges the two conjugacy classes when the rotation order is $2$ or $6$, because half that order is odd. In the square case, reflections along axes have split eigensublattices and reflections along diagonals have index two; both squared-length ratios are $1$. In the hexagonal case take $u=1$, $w=(1+i\sqrt3)/2$. Conjugation has invariant generator $u$ and anti-invariant generator $2w-u=i\sqrt3$; these span an index-two sublattice. Its ratio is $3$, and the other class has reciprocal ratio $1/3$. Finally $a=1$ in either model is square, so outside that case we can choose the representative $a>1$.
\end{proof}

\begin{corollary}[Refuted: the numerical guesses on 67/097--098]\label{A7-ratio-j}
With $Q$ the quadratic form used on 67/093 and 099, the hexagonal ratio is $3$, not $\sqrt3$. The usual modular invariant is
\[
 j(E_a)=j(i\sqrt a),\qquad
 j(E'_a)=j\left(\frac{1+i\sqrt a}{2}\right),
\]
and is unchanged by $a\mapsto1/a$ in either family. It is not $c\log a$ for either tentative nonzero constant in 67/097.
\end{corollary}
\begin{proof}
The quadratic form on the hexagonal basis has matrix
$\left(\begin{smallmatrix}1&1/2\\1/2&1\end{smallmatrix}\right)$, so $Q(2w-u)/Q(u)=3$. The parameter $\tau$ in each displayed period lattice gives the two formulas for $j$. Theorem~\ref{A7-real-classification} proves that passing to the reciprocal parameter changes the real structure on the same complex curve, hence leaves $j$ unchanged. On the other hand $\log(1/a)=-\log a$; in particular the square point has $a=1$ and $j=1728$, while $c\log1=0$. This disproves both tentative constants independently of their uncertain transcription. The exact quadratic calculation is in the attached program.
\end{proof}

\begin{theorem}[Proved: real points, torsors, and the topology on 67/101]\label{A7-real-topology}
The real locus and the first Galois cohomology group are
\[
 \begin{array}{c|c|c}
 E & E(\R) & H^1(\R,E)\\\hline
 E_a & (\R/\Z)\times(\Z/2\Z) & \Z/2\Z\\
 E'_a & \R/\Z & 0.
 \end{array}
\]
The quotients of the complex curves by their standard real structures are respectively an annulus and a M\"obius band. The unique nontrivial $E_a$-torsor has no real point; its antiholomorphic involution is $[z]\mapsto[\bar{z}+1/2]$, and its quotient is a Klein bottle. These torsors together with the two pointed families classify all smooth projective geometrically connected genus-one curves over $\R$.
\end{theorem}
\begin{proof}
In the split basis $u=1,v=i\sqrt a$, conjugation sends $(x,y)$ to $(x,-y)$. Thus fixed points modulo the lattice satisfy $2y\in\Z$, giving the two circles $y=0$ and $y=1/2$. In the nonsplit basis $u=1,w=(1+i\sqrt a)/2$, it sends $(x,y)$ to $(x+y,-y)$. The fixed-point condition is $y\in\Z$, leaving one circle.

Let $\Gamma=\operatorname{Gal}(\C/\R)$ and use the exact sequence of $\Gamma$-modules
\[
 0\longrightarrow\Lambda\longrightarrow\C\longrightarrow E(\C)
 \longrightarrow0.
\]
The higher cohomology of the real vector space $\C$ vanishes by averaging over the order-two group. The periodic resolution for a cyclic group therefore gives
\[
 H^1(\Gamma,E(\C))\simeq H^2(\Gamma,\Lambda)
 =\Lambda^\sigma/(1+\sigma)\Lambda.
\]
In the split case the numerator is $\Z u$ and the denominator $2\Z u$. In the nonsplit case $(1+\sigma)w=u$, so the quotient is zero. This is Galois cohomology of the elliptic curve, which classifies its torsors by \cite[IV, Theorem 7.8]{MilneEC2021A567}; alternatively the description follows directly by changing a chosen complex point on a torsor, changing the translation term of its real involution by a coboundary.

For the nontrivial split cocycle, $\tau[z]=[\bar{z}+1/2]$ satisfies $\tau^2[z]=[z+1]=[z]$. A fixed point would satisfy $z-\bar{z}=1/2+m+ni\sqrt a$ for integers $m,n$, impossible upon taking real parts. Thus the torsor has no real points. Conversely any real point on a torsor trivializes it by translation, so this nontrivial class is the only pointless form with that Jacobian.

The quotient by either reflection is a compact surface whose boundary is exactly its fixed locus: locally the quotient of complex conjugation is a half-disc. Its Euler characteristic is
\[
 \chi(E(\C)/\sigma)=\tfrac12\bigl(\chi(E(\C))+\chi(E(\R))\bigr)=0,
\]
as follows from a triangulation compatible with the fixed circles. The classification of compact surfaces therefore gives an annulus when there are two boundary circles and a M\"obius band when there is one. The free involution $\tau$ reverses orientation; its quotient is closed, nonorientable, and has Euler characteristic zero, hence is a Klein bottle. In particular the struck word ``sans'' on 67/101 must be restored: a nontrivial torsor gives a surface without boundary. Every genus-one real curve is a torsor under its Jacobian by \cite[IV, Theorem 7.19]{MilneEC2021A567}, and the preceding cohomology calculation exhausts all such torsors. Distinct parameters remain distinct because an isomorphism of real genus-one curves induces an isomorphism of their real Jacobians.
\end{proof}

The correspondence with Alling's classification in L9 is at the level of complete real isomorphism classes: the three species with geometrically connected complexification are annulus, M\"obius band, and Klein bottle, each with one positive real modulus. Alling \cite[preface and \S0.10]{Alling1981A567} explicitly states this classification and refers to the earlier Alling--Greenleaf treatment. The opened original 1971 preview defines dianalytic maps and proves their composition rule \cite[pp.~1--2, Proposition 1.1.3]{AllingGreenleaf1971A567}; its genus-one discussion, pp.~60--66, remains unopened. Our lattice proof above supplies the formulas and the uniqueness. The partial previews are distinguished from those uninspected classification pages.

\paragraph{The ranges of the real modular invariant.}
The two pointed families have ranges
\[
 \{j(E_a):a>0\}=[1728,\infty),\qquad
 \{j(E'_a):a>0\}=(-\infty,1728].
\]
Indeed a real short Weierstrass equation $y^2=x^3+Ax+B$ has
$j=1728\,4A^3/(4A^3+27B^2)$ and discriminant
$-16(4A^3+27B^2)$. A positive discriminant means three distinct real roots and two real components, and gives $j\ge1728$; a negative discriminant means one real root and one real component, and gives $j\le1728$. Conversely $A=-3$ with $0\le B^2<4$ realizes every $j\ge1728$; $A=-3$ with $B^2>4$ realizes every negative $j$; $A=0,B=1$ realizes $j=0$; and $A=3$ with $B^2\ge0$ realizes every $0<j\le1728$. The classification and real-locus calculation identify these with exactly the respective lattice families. The square endpoints are the two different real forms already listed.
